\section{Cotangent Message Passing}
\label{sec:cmp}

This section states the exact theorem proved by the current CMP construction.
It is a separator theorem for a linearized row space.

Let \(L_A,L_B,X_A,X_S,X_B\) be vector spaces over a field \(k\), and let
\begin{equation}
  J_A:L_A\to X_A,\quad
  J_{AS}:L_A\to X_S,\quad
  J_{BS}:L_B\to X_S,\quad
  J_B:L_B\to X_B
\end{equation}
be linear maps.  A pair \((\lambda,\mu)\in L_A\oplus L_B\) generates the
block-angular row combination
\begin{equation}
  \bigl(J_A\lambda,
  J_{AS}\lambda+J_{BS}\mu,
  J_B\mu\bigr).
  \label{eq:block-row-combination}
\end{equation}

\begin{definition}[Intrinsic CMP message]
\label{def:intrinsic-message}
The message from the \(A\)-side to the separator is
\begin{equation}
  \boxed{
  \cE(A\to S):=J_{AS}(\kernel J_A)\subseteq X_S.}
  \label{eq:intrinsic-message}
\end{equation}
Equivalently,
\begin{equation}
  \cE(A\to S)=
  \{J_{AS}\lambda:\lambda\in L_A,\ J_A\lambda=0\}.
\end{equation}
\end{definition}

\begin{theorem}[Linear separator factorization]
\label{thm:linear-separator-factorization}
For \(s\in X_S\) and \(b\in X_B\), the following are equivalent:
\begin{align}
 &\exists\lambda\in L_A,\ \exists\mu\in L_B:
 J_A\lambda=0,\quad
 s=J_{AS}\lambda+J_{BS}\mu,\quad b=J_B\mu;
 \label{eq:factor-left}\\
 &\exists e\in\cE(A\to S),\ \exists\mu\in L_B:
 s=e+J_{BS}\mu,\quad b=J_B\mu.
 \label{eq:factor-right}
\end{align}
\end{theorem}

\begin{proof}
From \cref{eq:factor-left}, set \(e=J_{AS}\lambda\); the condition
\(J_A\lambda=0\) implies \(e\in\cE(A\to S)\), yielding
\cref{eq:factor-right}.  Conversely, membership
\(e\in\cE(A\to S)\) supplies \(\lambda\in\kernel J_A\) with
\(e=J_{AS}\lambda\); substitution yields \cref{eq:factor-left}.
\end{proof}

The theorem has no rank, smoothness, or invertibility hypothesis.  It is exact
for the stated linear maps.

\subsection{Matrix form and basis independence}

Suppose row vectors of \(P_A^\perp\) form a basis of the relevant left null
space.  Then
\begin{equation}
  M_{A\to S}=P_A^\perp J_{AS}
  \label{eq:matrix-message}
\end{equation}
has row space \(\cE(A\to S)\).  A different null-space basis left-multiplies
\(P_A^\perp\) by an invertible matrix and therefore gives the same message
subspace.  Redundant rows may be removed without altering that subspace.

\begin{proposition}[Message dimension]
\label{prop:message-dimension}
If \(X_S\) is finite-dimensional, then
\begin{equation}
  \dim_k\cE(A\to S)\le\dim_kX_S.
\end{equation}
\end{proposition}

\begin{proof}
The message is a linear subspace of \(X_S\).
\end{proof}

\subsection{Recursive linearized elimination}

At a fixed point, each leaf sends a basis for the differential relations that
survive elimination of its internal coordinates.  At an internal node, child
messages are stacked with the node's own rows, internal columns are eliminated,
and the result is rank-reduced on the parent separator.  Induction using
\cref{thm:linear-separator-factorization} proves that the root message is exact
for the global \emph{linearized} row space.

For dense matrices of dimension at most \(k=P+K\), conventional null-space and
rank computations use \(O(k^3)\) field operations per bag.  With \(N\) bags,
this gives
\begin{equation}
  O\bigl(N(P+K)^3\bigr)
  \label{eq:cmp-arithmetic-count}
\end{equation}
field operations for the evaluated linear problem.  The coefficient field and
the cost of representing its elements remain part of any bit-complexity claim.

\subsection{Limit of the theorem}

The message in \cref{eq:intrinsic-message} is homogeneous first-order data.  It
contains directions generated by linear combinations of derivatives.  It does
not, by definition, contain the constant terms of the input polynomials or a
description of their real zero set.  The next section tests whether it can
nevertheless serve as a complete feasibility message and gives a negative
answer.
